\documentclass[11pt,a4paper]{article}

\usepackage[margin=1in]{geometry}
\usepackage[T1]{fontenc}
\usepackage{lmodern}
\usepackage{booktabs}
\usepackage{enumitem}
\usepackage{hyperref}
\usepackage{xcolor}
\usepackage{longtable}

\hypersetup{
  colorlinks=true,
  linkcolor=blue!50!black,
  urlcolor=blue!50!black,
  pdftitle={CompactifAI Aim: Implementation and Reproduction Report}
}

\title{CompactifAI Aim:\\Implementation and Reproduction Report}
\author{LLM Edge Compression Research}
\date{2026-11-06}

\begin{document}
\maketitle

\section{Aim}

Build an edge-efficient large-language-model compression pipeline that
replaces selected dense Transformer weight matrices with low-bond-dimension
Matrix Product Operator (MPO) representations, applies quantization to the
remaining layers, and recovers quality through calibration or fine-tuning.
The objective is a useful Pareto point: lower memory and storage, acceptable
target-device latency, and no unacceptable loss on the intended tasks.

\section{What the CompactifAI paper contributes}

\begin{enumerate}[leftmargin=*]
  \item \textbf{Layer tensorization:} reshape selected LLM matrices and
    decompose them into MPOs using sequential SVD truncation.
  \item \textbf{Bond dimension:} use the MPO bond dimension as the direct
    control for retained representational capacity and compression.
  \item \textbf{Layer sensitivity profiling:} protect sensitive layers or
    assign them larger bond dimensions.
  \item \textbf{Healing/retraining:} recover quality after local truncation by
    retraining the tensorized model.
  \item \textbf{Hybrid compression:} combine tensorization with quantization
    instead of forcing one method on every layer.
  \item \textbf{System evaluation:} report memory, parameters, training cost,
    speed, and quality together.
\end{enumerate}

The paper's reported LLaMA-2 memory reductions are paper results, not
guarantees for this repository. They must be reproduced with the same model,
data, hardware, and evaluation protocol before being claimed here.

\section{Current repository position}

The repository already provides:
\begin{itemize}[leftmargin=*]
  \item \texttt{ResearchMPOLinear}, which performs explicit TT-SVD with a
    maximum \texttt{bond\_dim};
  \item \texttt{ResearchMPOCompressor} for \texttt{nn.Linear} and GPT-2
    \texttt{Conv1D} modules;
  \item legacy rank-ratio MPO and low-rank baselines;
  \item mixed MPO plus dynamic int8 quantization;
  \item configurable protected-layer policies;
  \item bundle manifests and next-token evaluation support.
\end{itemize}

This is a reproduction scaffold rather than a complete CompactifAI
implementation. The principal missing capabilities are sensitivity-driven
rank selection, healing/fine-tuning, paper-style benchmark reporting, and
real-device performance validation.

\section{Implementation roadmap}

\subsection{Phase 0: Freeze a reproducible baseline}

Choose one legally accessible model and evaluation corpus. Record the model
revision, tokenizer, context length, software versions, hardware, random seed,
calibration data, evaluation data, batch size, sequence length, parameter
count, serialized size, peak memory, latency, throughput, and perplexity.

\textbf{Exit criterion:} the uncompressed model produces a saved baseline
report and can be rerun from a documented command.

\subsection{Phase 1: Validate MPO}

Use deterministic rectangular matrices representative of attention and MLP
projections. Test several bond dimensions, reconstruction error, forward-pass
equivalence, GPT-2 weight orientation, serialization, and reload behavior.

\textbf{Exit criterion:} error decreases as bond dimension increases, a
full-rank setting preserves outputs within numerical tolerance, and
save/reload is numerically equivalent.

\subsection{Phase 2: Map modules and set safe defaults}

Enumerate embeddings, attention projections, MLP projections, normalization,
and output-head modules for every model family. Begin conservatively by
protecting embeddings, the language-model head, the first block, and modules
that show a large quality drop. Produce a dry-run inventory of compressed and
skipped modules.

\subsection{Phase 3: Profile layer sensitivity}

Compress one layer or block at a time on a fixed calibration set. Record
quality change, reconstruction error, original parameters, compressed size,
and bond dimension. Use the resulting profile to allocate bond dimensions
under a memory budget.

\textbf{Exit criterion:} a sensitivity-aware policy gives better quality than
one global bond dimension at the same memory budget.

\subsection{Phase 4: Add healing}

Support both inexpensive calibration-only recovery and PEFT/QLoRA healing.
Keep the base model frozen where possible, record trainable parameters and
training cost, and evaluate on held-out data.

\textbf{Exit criterion:} healing recovers a measurable fraction of the quality
lost during compression without changing the evaluation protocol.

\subsection{Phase 5: Compare hybrid configurations}

\begin{enumerate}[leftmargin=*]
  \item original model;
  \item quantization only;
  \item MPO only;
  \item MPO plus quantization;
  \item MPO plus healing;
  \item MPO plus quantization plus healing.
\end{enumerate}

For each configuration report parameter count, serialized size, peak memory,
cold and warm latency, throughput, perplexity or next-token loss, downstream
scores, healing time, and trainable parameter count. Do not infer runtime
speed from parameter count; contraction overhead may dominate.

\subsection{Phase 6: Validate edge deployment}

Load the bundle in a fresh environment on the actual target device. Measure
end-to-end generation across context lengths, peak memory, sustained
throughput, warm latency, numerical stability, and unsupported-operation
fallbacks.

\section{Recommended experiment matrix}

\begin{longtable}{p{0.25\textwidth}p{0.32\textwidth}p{0.32\textwidth}}
\toprule
\textbf{Experiment} & \textbf{Settings} & \textbf{Question}\\
\midrule
MPO reconstruction & 2--4 sites and several bond dimensions &
Does the expected error/rank trade-off hold?\\
Global-rank baseline & One bond dimension for eligible layers &
How far can a simple policy go?\\
Sensitivity policy & Per-layer protection and bond dimensions &
Which layers dominate quality loss?\\
Quantization baseline & int8 and supported int4 &
What does quantization alone achieve?\\
Healing ablation & None, calibration, PEFT/QLoRA &
How much quality is recovered per training cost?\\
Hardware benchmark & CPU, CUDA, and target edge device &
Does lower storage reduce end-to-end cost?\\
Context sweep & Short, medium, and long prompts &
Does degradation grow with context length?\\
\bottomrule
\end{longtable}

\section{Proposed success criteria}

\begin{itemize}[leftmargin=*]
  \item At least 50\% lower serialized model size.
  \item At least 30\% lower peak inference memory.
  \item No more than one percentage point loss on the selected quality metric
    after healing.
  \item No unacceptable regression on selected downstream tasks.
  \item No warm-latency increase on the target device unless the aim is
    explicitly storage-only compression.
  \item Deterministic bundle loading and reproducible evaluation.
\end{itemize}

These are proposed engineering gates, not claims about the source paper.

\section{Deliverables}

\begin{enumerate}[leftmargin=*]
  \item A machine-readable sensitivity profile and rendered table.
  \item A baseline/compression benchmark JSON.
  \item Reproduction commands for each experiment.
  \item A healed compressed checkpoint and manifest.
  \item An edge-device benchmark log.
  \item A comparison against uncompressed and quantization-only baselines.
  \item Limitations covering licensing, data, unsupported layers, and hardware.
\end{enumerate}

\section{Immediate next actions}

\begin{enumerate}[leftmargin=*]
  \item Add MPO reconstruction, forward-equivalence, rectangular-matrix, and
    serialization tests.
  \item Add a dry-run module inventory to the compression CLI.
  \item Implement a one-layer sensitivity profiler with JSON output.
  \item Add perplexity/next-token-loss reporting.
  \item Implement a small calibration/healing experiment.
  \item Run the six-configuration matrix on one fixed model and target device.
  \item Publish benchmark files with the compressed bundle.
\end{enumerate}

\section{Conclusion}

Achieving the CompactifAI aim requires moving from ``an MPO decomposition
exists'' to ``a measured, sensitivity-aware, healed, and deployable workflow
exists.'' The current code provides the decomposition and hybrid-compression
foundation. The critical remaining work is profiling, recovery training,
reproducible evaluation, and real-device benchmarking, in that order.

\end{document}
